\subsection{Dataset Description}

The proposed RiskAdaptive framework is evaluated using the \texttt{TON-IoT}
network telemetry dataset, which contains diverse IoT network traffic
samples representing both normal and malicious activities. The dataset
includes multiple attack categories, enabling evaluation of the
framework under different threat scenarios.

The experiments consider attack detection capability as well as
risk-based prioritization performance across different attack types.


\subsection{Experimental Configuration}

The dataset is processed through preprocessing and feature
transformation stages before model training and evaluation. The
intrusion detection model is trained using labeled network traffic
samples, while the risk assessment module operates on the generated
detection outputs.

The evaluation pipeline consists of three stages:

\begin{enumerate}
\item Network traffic preprocessing and feature preparation.
\item Attack classification using machine learning based IDS.
\item Risk scoring and severity assignment using the RiskAdaptive
engine.
\end{enumerate}


\subsection{Baseline Comparison}

To evaluate the effectiveness of the proposed approach, comparisons
are performed against conventional machine learning based IDS models.
The baseline models include Logistic Regression and Random Forest
classifiers.

The comparison focuses on standard classification metrics including
accuracy, precision, recall, F1-score, false positive rate (FPR), and
false negative rate (FNR).


\subsection{Evaluation Metrics}

The following metrics are used for performance evaluation:

\begin{itemize}
\item Accuracy: Measures the overall correctness of predictions.
\item Precision: Represents the proportion of correctly identified
attacks among predicted attacks.
\item Recall: Measures the ability to detect actual attack instances.
\item F1-score: Represents the harmonic mean of precision and recall.
\item False Positive Rate: Measures incorrectly generated alerts.
\item False Negative Rate: Measures missed attack instances.
\end{itemize}


\subsection{Attack Scenario Evaluation}

In addition to overall classification performance, attack-specific
analysis is conducted for different threat categories including
backdoor, ransomware, and man-in-the-middle attacks.

This evaluation demonstrates the capability of RiskAdaptive to assign
different risk levels according to attack characteristics rather than
treating all detected events uniformly.
